%%%% Working TMLR revision. Default is the identified preprint.
% Pass \\def\\TMLRAnonymous{1} before input for the anonymous review draft.
\documentclass[10pt]{article}
\newif\ifanonymous
\ifdefined\TMLRAnonymous
  \anonymoustrue
  \usepackage{tmlr}
\else
  \anonymousfalse
  \usepackage[preprint]{tmlr}
\fi
\usepackage{url}
\usepackage[hidelinks]{hyperref}
\hypersetup{pdftitle={Response Geometry in Per-Step Integrated Gradients for Diffusion-Policy Vision-Language-Action Models}}
\ifanonymous
  \hypersetup{pdfauthor={}}
\else
  \hypersetup{pdfauthor={Arjun Bajpai}}
\fi
\usepackage[utf8]{inputenc}
\usepackage[small]{caption}
\usepackage{graphicx}
\usepackage{placeins}
\usepackage{amsmath}
\usepackage{amssymb}
\usepackage{amsthm}
\usepackage{booktabs}
\usepackage{algorithm}
\usepackage{algorithmic}
\usepackage{microtype}
\usepackage[table]{xcolor}

\renewcommand{\arraystretch}{1.15}

\urlstyle{same}

\title{Response Geometry in Per-Step Integrated Gradients\\for Diffusion-Policy Vision-Language-Action Models}

\author{\name Arjun Bajpai \email arjunbajpai2009@gmail.com \\
\addr Independent Researcher}

\begin{document}

\maketitle
% TMLR author separator and algorithmic boolean operator share this name.
\let\AND\relax

\noindent\textbf{Revision status (September 30, 2026).} This working draft\footnote{AI tools assisted with code review, analysis implementation, literature checks and drafting this revision. The author retains responsibility for all content.} corrects interpretations but retains labeled historical tables and figures pending canonical population reconciliation, real-model numerical validation, and a paired ranking/control study. It is not submission-ready.

\begin{abstract}
We study how scalar response choices affect the evaluation of Integrated Gradients (IG) for diffusion-policy vision-language-action models.
Our application to the Robotics Diffusion Transformer attributes a fixed-noise predicted action chunk to vision and language features and proprioceptive state at each policy call.
The primary scalar is an auxiliary quadratic action-discrepancy score, rather than an identified policy likelihood.
A retrospective analysis of saved perturbation curves separates the ranking used to select features from the response used to score their removal.
For the same recorded quadratic-target rankings, changing only the response to negative $\ell_2$ changes median deletion AUC from 0.4479 to 0.2908, reversing the historical threshold verdict without changing a selected feature.
Conversely, quadratic rescoring changes the recorded $\ell_2$ rankings' deletion AUC from 0.2830 to 0.4404.
An equal-feature example also shows that random-order normalized AUC need not be 0.5.
These results identify response geometry as a confound in cross-target attribution comparisons.
Historical RDT experiments additionally report nonzero action displacement, persistent correlations after backbone randomization, and target-dependent completeness residuals.
They do not isolate denoiser contraction, establish the dominance of input anchoring, or certify ranking efficacy without matched controls.
We distinguish reproducible saved-record analyses from historical summaries whose original inputs are unavailable, and identify numerical and provenance requirements for a controlled comparison.
The public artifact provides code and saved metrics for the analyses described here.
\end{abstract}

%% ================================================================
\section{Introduction}

Robotic manipulation policies that fuse vision, language, and proprioceptive state have advanced rapidly.
Vision-language-action (VLA) models such as RT-2~\citep{zitkovich2023rt2}, Octo~\citep{team2024octo}, and the Robotics Diffusion Transformer (RDT)~\citep{liu2024rdt} execute multi-step tasks from raw sensory input conditioned on natural-language instructions.
Diffusion-based policies~\citep{chi2023diffusion} have emerged as a particularly effective family within this class.
They generate action chunks through iterative denoising~\citep{ho2020ddpm} and can model multimodal action distributions; benchmark performance depends on the task and training setup.

Despite these advances, diffusion-policy VLAs remain opaque.
When a robot fails to grasp an object or misreads an instruction, diagnosing the contributions of vision, language and state is difficult from outcomes alone.
This limits deployment in safety-critical settings, hampers systematic debugging, and makes it harder for operators to decide when to intervene.

Post-hoc attribution offers a principled path toward interpretability without modifying the policy.
Integrated Gradients (IG)~\citep{sundararajan2017axiomatic} attributes a model's scalar output to its inputs by integrating gradients along a straight-line path from a baseline to the actual input.
It satisfies a completeness axiom that makes the attribution quantitatively checkable.
IG has been applied to image classifiers, question-answering systems and visual question answering models~\citep{sundararajan2017axiomatic,mudrakarta2018did}, motivating study of the corresponding construction in diffusion-policy episodic control.
Extending it to this setting raises three engineering obstacles.
The denoising chain samples noise, so the attributed scalar must specify how noise is conditioned on or averaged over.
The multimodal input space needs distinct baselines and interpolation strategies per modality.
Episodic control demands attribution at thousands of decision points across hundreds of episodes.

This paper addresses those engineering obstacles and studies the interpretation of attribution evaluation.
A perturbation curve depends on both the feature ranking and the scalar response to each perturbed input.
Changing the response from a quadratic discrepancy to a negative $\ell_2$ discrepancy can change normalized AUC even when the ranking and perturbed model outputs are fixed.
We demonstrate this directly on saved RDT curves and give an analytic counterexample to a universal random-ranking AUC of 0.5.
The question of whether a different attribution target produces a better ranking is separate and requires paired contexts, a common response, and informative controls.
Likewise, a small scalar discrepancy does not imply that the predicted action is unchanged, and a low input-shuffle correlation does not establish learned-model dependence.

Our contributions are as follows.
\begin{enumerate}
\item \textbf{Application.} A per-policy-call, per-modality IG construction for a fixed-noise diffusion sampler, with an explicit auxiliary action-discrepancy target and modality-specific feature baselines.
\item \textbf{Response-geometry analysis.} Exact within-ranking rescoring of recorded perturbation curves, showing that a historical deletion-AUC verdict can reverse without a change in feature ordering.
\item \textbf{Calibration counterexample.} An equal-contribution policy for which every ranking has continuous quadratic-response insertion and deletion AUC of $2/3$, establishing that 0.5 is not a universal chance reference.
\item \textbf{Empirical diagnostic record.} Historical completeness, perturbation, displacement, and randomization measurements, with explicit distinctions between reconstructible records, missing original artifacts, and unsupported causal interpretations.
\end{enumerate}
A controlled comparison of attribution rankings remains necessary to establish their efficacy beyond response geometry and simple ranking controls.

%% ================================================================
\section{Related Work}

\subsection{Attribution Methods}

Sundararajan et al.~\citep{sundararajan2017axiomatic} introduced IG with the sensitivity and implementation-invariance axioms and showed that straight-line path integration satisfies both.
Later work has addressed limitations of the vanilla method.
Jha et al.~\citep{jha2021smoother} computed attributions over neural stochastic differential equation models to produce smoother attribution maps, and later showed that shaping the injected noise further improves robustness~\citep{jha2022shaping}.
Walker et al.~\citep{walker2024idg} observed that gradients from saturation regions contribute noise and proposed Integrated Decision Gradients, which concentrates integration on the decision-relevant portion of the path.
For vision transformers, metric-driven attributions~\citep{walker2025metric} and information-flow decompositions~\citep{walker2025infoflow} improve patch-level explanations. Guided IG~\citep{kapishnikov2021guided} adapts the integration path to reduce pixel-level noise.
Petsiuk et al.~\citep{petsiuk2018rise} introduced the insertion and deletion curves we adopt as ranking-based faithfulness metrics.
The reliability of perturbation metrics is itself an active concern.
Hooker et al.~\citep{hooker2019roar} evaluate attributions by retraining on ablated inputs.
Rong et al.~\citep{rong2022road} show that deletion without retraining leaks information through the shape of the mask itself, and Tomsett et al.~\citep{tomsett2020sanity} find that common saliency metrics are themselves unreliable.
Our structural findings add a diffusion-policy instance of that caution, where the readout convention rather than the ranking drives the magnitude verdict.
Adebayo et al.~\citep{adebayo2018sanity} proposed model and data randomization tests to assess whether saliency maps depend on learned parameters and training data.
Their results caution against interpreting visually plausible maps as evidence of that dependence.
Our backbone randomizations are limited tests because learned encoders and adaptors remain in place.
For diffusion generative models in the image domain, cross-attention maps relate generated pixels to prompt tokens~\citep{tang2023daam}.
This differs from attributing a scalar discrepancy of a predicted control chunk.

These methods target classifiers, language models, or image generators.
We extend IG to diffusion-policy VLAs, which requires handling stochastic denoising, action-chunk outputs, and multimodal inputs inside an episodic control loop.

\subsection{Vision-Language-Action Models}

VLA models unify perception, language understanding, and action generation.
RT-2~\citep{zitkovich2023rt2} fine-tunes a vision-language model to emit discretized actions as token sequences.
Octo~\citep{team2024octo} trains a generalist policy on a large multi-robot dataset with a transformer backbone.
Diffusion Policy~\citep{chi2023diffusion} frames action generation as conditional denoising diffusion~\citep{ho2020ddpm}.
It produces action chunks by iterative refinement rather than single-step regression.
RDT~\citep{liu2024rdt} scales this approach with a Diffusion Transformer (DiT) backbone, a SigLIP vision encoder~\citep{zhai2023siglip}, and a T5-XXL language encoder~\citep{raffel2020t5}.
RDT processes six camera slots (three camera views over a two-frame history) through SigLIP into 4{,}374 vision tokens.
It combines them with up to 1{,}024 language tokens and a 128-dimensional unified state vector, and it generates 64-step action chunks through a five-step DPM-Solver++~\citep{lu2022dpm} denoising schedule.
This multimodal, multi-step architecture makes RDT a compelling but challenging target for attribution.

\subsection{Interpretability in Robotics}

Prior work on explaining robotic policies has focused on reinforcement learning agents in discrete-action settings.
Greydanus et al.~\citep{greydanus2018visualizing} applied perturbation-based saliency to Atari agents, and counterfactual methods~\citep{amitai2024explaining} contrast the outcomes of executed and alternative actions for policy explanation.
In continuous control, VisualBackProp~\citep{bojarski2018visualbackprop} visualizes convolutional neural network (CNN) saliency for end-to-end driving policies.
That work addresses convolutional driving models rather than a multimodal diffusion sampler.
Recent VLA studies examine different explanation targets.
Zhang et al.~\citep{zhang2026embodied} estimate visual influence using randomized perturbations and action-discrepancy responses, with explicit assumptions for relating mean-squared error to distributional divergence.
VLA-Trace~\citep{shi2026vlatrace} combines representation comparisons, attention interventions and behavioral probes.
IMPACT-VLA~\citep{kim2026impact} attributes phase-modality interventions to downstream task success through counterfactual closed-loop execution.
Our estimand is the discrepancy of a fixed-noise predicted action chunk at a fixed context, not an explanation of eventual task success.
The present study examines how response geometry affects the evaluation of IG rankings for this estimand; it does not claim to originate VLA interpretability.

%% ================================================================
\section{Method}

\subsection{Integrated Gradients Preliminaries}

Given a differentiable function $F \!:\! \mathbb{R}^n \!\to\! \mathbb{R}$, an input $x \in \mathbb{R}^n$, and a baseline $x' \in \mathbb{R}^n$, the integrated gradient along dimension $i$ is
\begin{equation}
\mathrm{IG}_i(x, x') = (x_i - x'_i) \int_0^1 \frac{\partial F}{\partial x_i}\bigg|_{x' + \alpha(x - x')} d\alpha .
\label{eq:ig}
\end{equation}
The \emph{completeness axiom} guarantees $\sum_i \mathrm{IG}_i(x, x') = F(x) - F(x')$.
We approximate the integral with a uniform Riemann sum of $m$ steps, giving
\begin{equation}
\mathrm{IG}_i \approx \frac{x_i - x'_i}{m+1} \sum_{r=0}^{m} \frac{\partial F}{\partial x_i}\bigg|_{x' + \frac{r}{m}(x - x')} .
\label{eq:riemann}
\end{equation}
This equal-weight endpoint average describes the historical implementation.
Convergence requires regularity of the path integrand; finite-precision accumulation introduces additional error that increasing $m$ does not isolate or necessarily reduce.

\subsection{A Deterministic Target Through the Denoising Chain}
\label{sec:target}

At policy call $t$, write the sampler output as $g_\theta(o_t,z_t)$ for observation $o_t$ and fixed diffusion noise $z_t$.
We denote this predicted action chunk by $\mu_t$ for consistency with historical notation; it is a sample conditioned on fixed noise, not an identified conditional mean.
The reference chunk $a_t=g_\theta(o_t,z_t)$ is computed once without gradients.
For an intervened input, the primary attribution target is the auxiliary score
\begin{equation}
Q_t=-\frac{1}{2\sigma^2D}\lVert a_t-\mu_t\rVert^2,
\label{eq:target}
\end{equation}
with $\sigma^2=1$ and $D=512$ active entries: 64 predicted time steps over eight active dimensions.
This is the nonconstant part of a Gaussian log kernel averaged over entries, with its log-normalizing constant omitted.
It has not been derived as the diffusion policy's conditional log likelihood.
Historical code and files call this target \texttt{logpi}; we use $\Delta Q$ for its change.
Attribution concerns the complete predicted chunk, of which the execution loop applies a subsampled prefix before replanning or termination.

We hold the noise realization fixed across the interpolation path to evaluate one deterministic scalar function.
Attribution of an expectation over noise would be a different estimand and requires a justified gradient-expectation interchange and Monte Carlo treatment.
Independent noise at each path point does not compute the fixed-noise quantity considered here.
The historical implementation replaced nonfinite gradient entries by zero.
That operation can corrupt individual coordinates while cancellation leaves the scalar completeness residual small or even zero; completeness is not a certificate of coordinate accuracy.

\subsection{Alternative Scalar Targets}
\label{sec:targets}

We compare the quadratic target with two alternative action-discrepancy targets,
\begin{equation}
y_t^{\ell_2}=-\sqrt{\lVert a_t-\mu_t\rVert_2^2+\epsilon},
\qquad y_t^{\mathrm{max}}=-\max_j|a_{t,j}-\mu_{t,j}|,
\end{equation}
where $\epsilon=10^{-12}$ is the implemented stabilizer and both targets use the same $D$ active entries.
The stabilized $\ell_2$ target is smooth in the residual and has zero gradient at exact agreement; away from the stabilizer scale its action-space gradient approaches unit norm.
The maximum-deviation target, abbreviated maxdev, uses a different norm and is not a monotone transform of $\ell_2$ in general: residuals $(1,1)$ and $(1.1,0)$ reverse order.
The quadratic and stabilized $\ell_2$ responses are algebraically related for a fixed perturbation, although their gradients can induce different attribution rankings.
Their units and scales differ, so a larger magnitude alone does not demonstrate greater information or operational usefulness.
Section~\ref{sec:targets-results} evaluates all three.

\subsection{Per-Step Episode Attribution}

RDT outputs action chunks of 64 steps, subsampled $4\times$ to 16 environment steps.
We run IG once per \emph{policy call}, not per environment step, because the chunk is the atomic unit of decision.
Per-call attribution costs about one minute per decision on RDT-170M at $m{=}64$.
Attributing every environment step would multiply that cost by the 16 subsampled steps per chunk and put a full pass beyond practical budgets.
The per-call choice is what makes full-episode evaluation tractable.

At each policy call we compute three independent attributions through modality-specific forward functions.
\begin{itemize}
\item \textbf{Vision.} Freeze language and state and interpolate the adapted image tokens from the baseline to the real observation.
\item \textbf{Language.} Freeze vision and state and interpolate the adapted language tokens from a minimal two-token baseline encoded through T5-XXL to the real instruction embedding.
\item \textbf{State.} Freeze vision and language and interpolate the 128-dimensional unified state vector (eight active dimensions for the Franka arm, namely seven joints plus the gripper) from zeros to its real value. The state-adaptor multilayer perceptron (MLP) sits inside the forward function to preserve per-joint structure.
\end{itemize}
Algorithm~\ref{alg:perstep} summarizes the procedure.

\begin{algorithm}[t]
\caption{Per-step IG for diffusion-policy VLAs.}
\label{alg:perstep}
\textbf{Input.} Policy $\pi_\theta$, environment $\mathcal{E}$, integration steps $m$, target $y \in \{Q, \ell_2, \mathrm{maxdev}\}$\\
\textbf{Output.} Attributions $\{A_t^V, A_t^L, A_t^S\}$ at every policy call.
\begin{algorithmic}[1]
\STATE $o_0 \leftarrow \mathcal{E}.\text{reset}()$
\FOR{policy call $t = 0, 1, \ldots$}
    \STATE Encode $o_t$ into vision $v_t$, language $l_t$, state $s_t$
    \STATE Compute $a_t \leftarrow \pi_\theta(o_t)$ with fixed seed, no grad
    \STATE Define $F$ from $a_t$, the fixed-noise predicted chunk $\mu_t$, and the chosen target
    \FOR{modality $M \in \{V, L, S\}$}
        \STATE Set baseline $x'_M$, input $x_M$
        \STATE Freeze the other two modalities
        \STATE $A_t^M \leftarrow \text{IG}(F_M, x_M, x'_M, m)$ \hfill Eq.~\ref{eq:riemann}
    \ENDFOR
    \STATE Record $A_t^V, A_t^L, A_t^S$ and completeness errors
    \STATE Apply action chunk and advance $\mathcal{E}$ by 16 steps
\ENDFOR
\end{algorithmic}
\end{algorithm}

\paragraph{Baseline selection.}
The vision baseline is a uniform gray image in the external-camera slot only.
The remaining five camera slots carry constant background images and are frozen at their real values, so IG attributes only to the informative slot.
The language baseline is a minimal two-token sequence (the tokenizer's padding and end-of-sequence tokens, since T5 has no beginning-of-sequence token) encoded through T5-XXL rather than a zero vector.
The encoder output is zero-padded to the stored embedding length, and the language forward passes attend over the task's real token span.
The attended baseline is therefore the encoded token pair followed by zero vectors at the remaining attended positions.
A zero vector is numerically permitted by the GELU-based language adaptor but defines a different representation baseline; preliminary efficacy comparisons are not independently reconstructible.
The state baseline is a zero vector in the normalized joint space, a numerically safe choice that keeps per-joint attributions readable.
Section~\ref{sec:anchoring} quantifies sensitivity to the vision baseline choice.

\subsection{Evaluation Protocol}
\label{sec:metrics}

We evaluate along three axes against fixed quantitative bars.
The historical numeric bars are study-specific descriptive targets, not community standards or calibrated hypothesis tests.
Random ranking does not generally yield AUC 0.5 after endpoint normalization.
For an equal-feature additive scalar policy with all-one input, zero baseline, and reference action one, deletion fraction $q$ gives normalized quadratic response $1-q^2$, while insertion gives $1-(1-q)^2$.
Every ordering has continuous AUC $2/3$ for both curves; on the historical token grid the values are about 0.6597, exceeding the historical vision insertion bar without ranking information.
Reported 95\% intervals use episode-resampled percentile bootstrap intervals with $B=10{,}000$ and a fixed seed; the episode identities must correspond to independent sampled contexts for the clustering interpretation to hold.

\paragraph{Correctness.}
We measure the completeness error
\begin{equation}
\text{Err} = \frac{\left| \sum_i \mathrm{IG}_i - (F(x) - F(x')) \right|}{\left| F(x) - F(x') \right|}
\end{equation}
per modality at every policy call.
Completeness is evaluated per modality, with the other two modalities frozen, rather than as one joint sum over all modalities, because each modality has its own forward function and baseline.
The bars are a median error at or under 3\% and at least 90\% of cases at or under 3\%.

\paragraph{Faithfulness.}
We run two perturbation tests.
\emph{Top-$k$ deletion} ranks inputs by $|\mathrm{IG}|$, replaces the top $k\%$ with the baseline for $k \in \{1, 5, 10\}$, and measures $\Delta F = F(\text{deleted}) - F(\text{original})$, which is nonpositive for a self-referenced discrepancy in exact arithmetic and is not itself evidence of ranking quality.
The historical bar at $k{=}5\%$ is $\Delta Q \leq -0.5$ under the auxiliary per-entry Gaussian convention.
\emph{Insertion/deletion AUC}~\citep{petsiuk2018rise} progressively inserts tokens starting from the baseline, or deletes them starting from the input, in IG-ranked order over $k \in \{0, 1, 5, 10, 20, 30, 50, 75, 100\}\%$.
We compute the AUC after rescaling $F$ so that the full input maps to 1 and the baseline to 0.
The ranked units are the 729 external-camera patch tokens for vision and the real instruction tokens for language.
The bars are vision insertion at or above 0.60 and deletion at or below 0.40, and language insertion at or above 0.55 and deletion at or below 0.45.
The auxiliary score scale of $\Delta Q$ is set by the target's normalization conventions, namely $\sigma^2{=}1$ and the per-entry mean over $D$.
An alternative sum-form convention without the $1/D$ factor would scale every reported $\Delta Q$ by $D{=}512$ (Section~\ref{sec:contraction}).
The historical magnitude bar applies only to the quadratic target under that convention, not to alternative-target units.

\paragraph{Sanity.}
\emph{C1 (model randomization)} re-initializes model weights with Xavier-uniform draws and re-runs IG.
Attributions should decorrelate from the original, with Spearman $\rho \leq 0.2$.
We run three randomization variants.
The standard variant re-initializes the output projection of the final action layer and recomputes the target under the perturbed model.
The frozen-target variant re-initializes the same projection but keeps the original model's reference action.
The cascade variant re-initializes every backbone weight matrix and positional embedding (107 tensors) with the frozen target.
\emph{C2 (input randomization)} applies a deterministic pixel permutation to the external-camera image and a post-encoding token-order shuffle to the language embedding, then re-runs IG with the same $\rho \leq 0.2$ bar.
C1 follows the model-randomization test of \citeauthor{adebayo2018sanity}~\citep{adebayo2018sanity}, while C2 is an inference-time input-randomization check in the same spirit.
Correlations are computed over the 729 external-camera patch scores for vision, the real instruction tokens for language, and the eight active state dimensions.
All sanity re-runs compute the perturbed attribution at $m{=}16$ on the first 50 steps of each task-seed base run and correlate it against that run's original $m{=}64$ attribution.
The standard C1 and C2 rows use the main pass as the base run, while the frozen and cascade variants use the regeneration run described in Section~\ref{sec:setup}.
The standard C1 result jointly changes weights and integration budget; it does not isolate the effect of changing $m$ alone.
The stored attribution sidecars, the per-policy-call attribution tensors, reference action, observation, and proprioceptive state saved alongside each run, encoded neither the integration budget nor the episode seed in their filenames until a tagging fix that postdates every run described in this paragraph.
Both consequences are disclosed.
First, the $m{=}128$ extension ran concurrently with the variant sanity stages on PegInsertionSide and PickSingleYCB, so some frozen-target and cascade rows for those two tasks may correlate against $m{=}128$ attributions rather than $m{=}64$ ones.
Second, the two seeds of each pre-fix run wrote to the same sidecar paths, with the second seed writing last, so every pre-fix stage that loaded sidecars for the first seed is exposed to having consumed the second seed's tensors, mixing the second seed's attribution rankings and observations with the first seed's noise seeding and row filtering.
A released consistency fingerprint compares each step record's completeness gap with the faithfulness record's recomputed baseline value.
The seed-42 $m=128$ faithfulness records disagree with their own step records on all tested rows, whereas post-fix records agree.
A discrepancy reveals an incompatibility; scalar equality alone cannot authenticate an exact observation, checkpoint, or noise identity.
Subset sensitivity analyses leave historical threshold descriptions unchanged, but do not restore missing identities or validate independent-cluster inference for reused contexts.

\paragraph{Baseline sensitivity.}
For the vision modality we recompute IG under three baselines per step (black, gray, and a Gaussian blur of the actual observation) and report pairwise Spearman correlations of the patch rankings.

%% ================================================================
\section{Experimental Setup}
\label{sec:setup}

\paragraph{Models.}
RDT-170M~\citep{liu2024rdt} (14 DiT layers, 1024-dim hidden) is the primary model, used with pretrained-only weights.
This is deliberate.
A policy without task-specific fine-tuning exercises the attribution machinery on realistic failure behavior, and its near-zero task success (1 success in 400 episodes) supplies the failure side of the qualitative analysis.
RDT-1B (28 layers, 2048-dim) serves as the scale check on PickCube with the original authors' fine-tuned weights, a checkpoint that can solve the task.
An initial scale-check run produced 7/20 successes.
A re-run of the same checkpoint for the faithfulness evaluation, on the follow-up machine and under the deviations recorded below, produced 9/20.
A Month 5 reproduction run of the same checkpoint and seeds with real T5-XXL embeddings on a third machine produced 8/20 successes, and its records are committed.
Its vision faithfulness medians fall inside the episode bootstrap intervals of the committed run (insertion 0.919 against 0.926 and deletion 0.264 against 0.248).

\paragraph{Tasks.}
We evaluate on four ManiSkill3~\citep{tao2024maniskill3} manipulation tasks with a simulated Franka Panda seven-degree-of-freedom (7-DoF) arm.
PickCube-v1 grasps a cube and lifts it to a target.
StackCube-v1 stacks a red cube on a green cube.
PegInsertionSide-v1 inserts a peg into a hole from the side.
PickSingleYCB-v1 grasps a Yale-CMU-Berkeley (YCB) object.
Each task is conditioned on a natural-language instruction encoded offline through T5-XXL.
Tables abbreviate PegInsertionSide-v1 and PickSingleYCB-v1 as PegInsertion and PickYCB.
We record five protocol deviations up front.
The run campaigns are labeled by project month, Month 4 through Month 7, and the released records and data README use the same labels.
The original task list named OpenDrawer, which exists only as a Fetch-mobile-manipulator variant in the installed ManiSkill3, so it was dropped.
PickSingle does not exist under that name, so the published PickSingleYCB-v1 was substituted.
The recorded TurnFaucet attempt used an incorrect controller, and no corrected rerun is documented.
Its low norms and unsuccessful outcomes therefore do not establish task incompatibility or justify a capability-based exclusion; the four-task selection is a historical scope limitation.
The Month 4 follow-up machine could not host T5-XXL, so every follow-up run on it loaded the authors' precomputed instruction embeddings and a zeroed placeholder in place of the encoded two-token language baseline.
The affected language numbers are the seed 42 RDT-1B language values reported with Table~\ref{tab:faithfulness}, the variant language correlations in Table~\ref{tab:sanity}, and the regeneration-run and $m{=}128$-overlay language distributions in Figure~\ref{fig:completeness}.
Comparisons between them and main-pass language numbers therefore span two baselines, while the vision and state baselines are unchanged.
The Month 6 runs at seeds 142 and 242 encoded their instruction embeddings and the two-token language baseline with T5-XXL on their own machine, so the pooled 1B language row in Table~\ref{tab:faithfulness} spans both language-baseline provenances.
The follow-up PickSingleYCB runs also reused the PickCube instruction embedding because the task's own embedding was unavailable, so no per-task language metrics are tabulated from them.
The PickSingleYCB language distributions in Figure~\ref{fig:completeness} carry that proxy embedding, and the pooled variant language correlations include those proxy rows.

\paragraph{Protocol.}
The full evaluation pass is 4 tasks $\times$ 2 seeds (42, 142) $\times$ 50 episodes at $m{=}64$.
It yields roughly 1{,}900 attributed policy calls, of which 1{,}341 are action-norm-selected and enter the quantitative tables.
The historical action-norm filter retains reference norms at least 15; we call these action-norm-selected rows.
The norm of one action does not bound the completeness denominator $|F(x)-F(x')|$, so this filter is not a conditioning or signal guarantee.
Seven follow-up runs extended the protocol.
An $m{=}128$ subset covered PickCube and StackCube (2 seeds $\times$ 8 episodes) and was later extended to PegInsertionSide and PickSingleYCB at the same scope.
A configuration fix raising the episode-step ceiling landed between the $m{=}64$ pass and all later runs.
As a result, later episodes typically run to 25 policy calls while the original $m{=}64$ episodes were shorter.
Depth-truncated comparisons in Section~\ref{sec:budget} truncate the $m{=}128$ rows to the $m{=}64$ episode shape.
A regeneration run repeated all four tasks at $m{=}64$ (2 seeds $\times$ 15 episodes, post-fix configuration) to supply base attributions for the variant sanity checks, the baseline study, and the target ablation.
The target ablation re-ran PickCube at the same scope once per alternative target.
The 1B faithfulness run covered PickCube at three evaluation seeds (42, 142, 242), 20 episodes each at $m{=}64$, with the seed 42 row kept from the earlier follow-up run and seeds 142 and 242 added later.
A displacement run measured the action change under top-ranked deletion directly in action space on PickCube and StackCube at both model scales.
It swept solver steps $T \in \{1, 2, 3, 5, 10, 20\}$, the denoising-step count of Section~\ref{sec:target}, and deletion fractions up to 20\% on vision and language, with rankings taken from the $T{=}5$ attributions.
The reported vision curves pool the 1B runs over seeds 42 and 142, use seed 42 at 170M, and carry 8 to 40 episodes per curve.
A Month 5 reproduction run repeated the 1B faithfulness with real T5-XXL embeddings.
A Month 7 verification pass repeated the full evaluation pass, the faithfulness stage, and the standard sanity checks at the main-pass scope with seed-tagged sidecars, together with a clean rerun of the $m{=}128$ seed-42 arm, following the external audit of the sidecar filename scheme disclosed in Section~\ref{sec:metrics}, and its records are committed.

\paragraph{Hardware.}
The main pass ran on a single RTX 4090 (24\,GB).
The original $m{=}128$ subset (PickCube and StackCube) ran on an RTX 5090 (32\,GB, PyTorch 2.8 with CUDA 12.8).
The $m{=}128$ extension to the other two tasks and the remaining Month 4 follow-up experiments ran on a machine with two RTX 4090s.
The displacement measurement and the Month 5 reproduction run used a RunPod RTX PRO 6000 Blackwell (96\,GB), which hosted T5-XXL and therefore used real instruction embeddings.
The Month 6 runs (the seed 142 and 242 RDT-1B faithfulness runs, the nominally matched $Q$ ablation, and the seed 142 displacement runs) ran on a second RunPod RTX PRO 6000 pod.
That pod encoded its own T5-XXL instruction embeddings and two-token language baseline.
T5-XXL instruction embeddings are precomputed offline, so the 4.8B-parameter encoder never resides in GPU memory during attribution.

%% ================================================================
\section{Evaluation Results}

This section retains historical summaries while distinguishing available records from missing original data.
Passing a study-specific bar does not by itself establish attribution efficacy.

\subsection{Correctness}

Table~\ref{tab:correctness} preserves the original main-pass completeness summary; its raw records are unavailable, so it is not primary reconstructible evidence.
Figure~\ref{fig:completeness} uses a separate committed regeneration campaign with its own population and baseline deviations.
Historical vision and language medians range from 1.87 to 2.54\%, while state medians range from 4.53 to 7.46\%.
The 90\%-of-cases bar is missed at $m=64$.
These scalar residuals do not establish coordinate accuracy or identify causes of numerical error.
Historical gradients accumulated in the input dtype, and nonfinite substitutions were not counted in saved records.
Same-context numerical validation is needed for budget-only conclusions and current map validation.
\begin{table}[tbp]
\centering
\footnotesize
\setlength{\tabcolsep}{3.5pt}
\begin{tabular}{@{}l r rr r rr r rr@{}}
\toprule
 & & \multicolumn{2}{c}{\textbf{Vision}} & & \multicolumn{2}{c}{\textbf{Language}} & & \multicolumn{2}{c}{\textbf{State}} \\
\cmidrule(lr){3-4} \cmidrule(lr){6-7} \cmidrule(lr){9-10}
Task & $n$ & med & $\leq$3\% && med & $\leq$3\% && med & $\leq$3\% \\
\midrule
PegInsertion & 425 & 2.11 & 75.8 && 1.95 & 73.9 && 4.70 & 37.9 \\
\rowcolor{gray!6}
PickCube & 264 & 2.02 & 68.6 && 2.08 & 67.0 && 5.12 & 34.8 \\
PickYCB & 348 & 2.06 & 69.8 && 2.06 & 67.8 && 4.53 & 36.5 \\
\rowcolor{gray!6}
StackCube & 304 & 2.54 & 59.9 && 1.87 & 75.0 && 7.46 & 25.3 \\
\bottomrule
\end{tabular}
\caption{Completeness error (\%) at $m{=}64$ on RDT-170M, action-norm-selected rows ($|$ref action$| \geq 15$). The ``med'' column is the median error. The ``$\leq$3\%'' column is the fraction of rows at or under the 3\% bar.}
\label{tab:correctness}
\end{table}

\begin{figure}[tbp]
\centering
\includegraphics[width=0.85\linewidth,height=0.38\textheight,keepaspectratio]{figures/fig_a_completeness.png}
\caption{Completeness error distributions by task and modality over all policy calls of the committed regeneration run at $m{=}64$ (2 seeds $\times$ 15 episodes per task). The PegInsertionSide and PickSingleYCB panels carry $m{=}128$ overlays, and the PickCube panel carries the seed 42 RDT-1B scale-check rows. Vision and language mass concentrates at low error with tails extending past the 3\% line. State shows a longer right tail. The PickSingleYCB language histograms reflect the proxy PickCube instruction embedding recorded in Section~\ref{sec:setup}. All language distributions in this figure come from runs that used the zeroed language baseline recorded in Section~\ref{sec:setup}. Errors are clipped at 0.5 (50\%) for display, so the rightmost bin pools all larger errors.}
\label{fig:completeness}
\end{figure}

\subsection{Faithfulness}
\label{sec:faithfulness}

Table~\ref{tab:faithfulness} reports faithfulness at $m{=}64$ on the same action-norm-selected population, with the RDT-1B scale check in the bottom block.
The 170M block of Table~\ref{tab:faithfulness} is a historical summary whose original rows are unavailable.
Vision insertion AUC ranges from 0.792 to 0.866 and deletion AUC from 0.290 to 0.345.
These values cross the historical bars, but no matched random-ranking comparison establishes an advantage over a null baseline.
The interventions replace adapted patch tokens, not pixels.
Language insertion ranges from 0.344 to 0.504.
With only 16 to 30 instruction tokens, requested percentage interventions differ from realized fractions.
Neither token count nor the cross-checkpoint comparison isolates a cause; the latter also spans baseline and population changes.

\begin{table}[tbp]
\centering
\footnotesize
\setlength{\tabcolsep}{4pt}
\begin{tabular}{@{}ll r r rr@{}}
\toprule
\textbf{Task} & \textbf{Mod.} & $\boldsymbol{n}$ & $\boldsymbol{\Delta Q}$ & \textbf{Ins} & \textbf{Del} \\
\midrule
PegInsertion & Vision & 425 & $-$0.0018 & \textbf{0.792} & \textbf{0.290} \\
& Lang. & 425 & $-$0.0039 & 0.344 & \textbf{0.249} \\
\midrule
\rowcolor{gray!6}
PickCube & Vision & 264 & $-$0.0032 & \textbf{0.856} & \textbf{0.317} \\
\rowcolor{gray!6}
& Lang. & 264 & $-$0.0004 & 0.504 & 0.477 \\
\midrule
PickYCB & Vision & 348 & $-$0.0013 & \textbf{0.843} & \textbf{0.345} \\
& Lang. & 348 & $-$0.0003 & 0.368 & 0.477 \\
\midrule
\rowcolor{gray!6}
StackCube & Vision & 304 & $-$0.0039 & \textbf{0.866} & \textbf{0.336} \\
\rowcolor{gray!6}
& Lang. & 304 & $-$0.0016 & 0.378 & 0.539 \\
\midrule
\midrule
RDT-1B s42 & Vision & 320 & $-$0.00016 & \textbf{0.926} & \textbf{0.248} \\
RDT-1B s142 & Vision & 480 & $-$0.00012 & \textbf{0.937} & \textbf{0.293} \\
RDT-1B s242 & Vision & 399 & $-$0.00018 & \textbf{0.954} & \textbf{0.314} \\
\rowcolor{gray!6}
RDT-1B pool & Vision & 1{,}199 & $-$0.00015 & \textbf{0.940} & \textbf{0.288} \\
\rowcolor{gray!6}
RDT-1B pool & Lang. & 1{,}199 & $-$0.00012 & 0.541 & \textbf{0.267} \\
\bottomrule
\end{tabular}
\caption{Faithfulness at $m{=}64$. $\Delta Q$ is the median target change after top-5\% deletion, in auxiliary per-entry log-kernel units. Bold AUC values cross their historical bars (vision insertion $\geq$0.60 and deletion $\leq$0.40, language $\geq$0.55 and $\leq$0.45). The top block is RDT-170M with both seeds on action-norm-selected rows ($|$ref action$| \geq 15$). The bottom block is the RDT-1B scale check with the authors' fine-tuned weights on three evaluation seeds (42, 142, 242, with 9, 1, and 5 of 20 successes), on all policy-call rows per seed. The historical filter was not applied at 1B because its action norms sit below 15 throughout, creating a different selected population. The seed 42 language values were computed against the zeroed language baseline recorded in Section~\ref{sec:setup}, while the seed 142 and 242 runs encoded their instruction embeddings and the two-token baseline with T5-XXL on their own machine. The pooled language row therefore spans the two language-baseline provenances. The vision rows are unaffected by that difference. Each of the three vision rows crosses both historical bars; this is not a test of cross-checkpoint robustness or efficacy against a null ranking. The pooled vision 95\% episode bootstrap interval is 0.926 to 0.954 for insertion and 0.250 to 0.312 for deletion. Language insertion is seed-sensitive. It is 0.562, 0.583, and 0.519 at the three seeds, so it passes the 0.55 bar at two of the three seeds. The pooled value 0.541 sits just under the bar with a 95\% interval of 0.449 to 0.605 that straddles it. The committed Month 5 reproduction of seed 42 with real T5-XXL embeddings records a language insertion median of 0.515, below the bar, so the two of three count is tied to the language-baseline provenance rather than to the seeds alone. We therefore report the language insertion pass as seed-sensitive rather than robust, while the vision verdicts and the language deletion pass hold at every seed. The RDT-170M block predates the sidecar seed-tagging fix of Section~\ref{sec:metrics}, and its seed-42-labeled halves are presumed to have consumed seed-142 sidecars, mixing that seed's rankings and observations with seed-42 noise seeding and signal filtering. The released verification records from the full protocol rerun with seed-tagged sidecars (Section~\ref{sec:setup}) reproduce every vision verdict in this block, with insertion from 0.766 to 0.868 and deletion from 0.286 to 0.362 across the four tasks, and reproduce the language misses. One borderline cell crosses its bar: PickCube language deletion is 0.438 in the verification pass against the tabulated 0.477, moving to the passing side of the 0.45 bar, consistent with reading the 170M language metrics as mixed. These are evaluation seeds on one checkpoint, so they measure rollout and sampling variance, not generalization across training initializations.}
\label{tab:faithfulness}
\end{table}

Historical top-5\% quadratic-score changes have magnitudes around $10^{-3}$ or smaller.
Pooled 1B values are smaller still, despite recorded successful episodes and changes in predicted actions.
This observational contrast changes checkpoints, training, action regimes and filtering rules; it is not a controlled competence effect.
Magnitudes describe the chosen score and require a stated convention and operational criterion.

\begin{figure}[tbp]
\centering
\includegraphics[width=0.85\linewidth,height=0.38\textheight,keepaspectratio]{figures/fig_c_auc_curves.png}
\caption{Normalized insertion (higher is better) and deletion (lower is better) curves over the token-fraction grid, aggregated per modality over all faithfulness rows, including the $m{=}128$ subset, which is a broader population than Table~\ref{tab:faithfulness}. The legend's $m{=}64$ label comes from the file tagging in the released notebook and does not describe the pooled population.}
\label{fig:auc}
\end{figure}

\subsection{Sanity Checks}
\label{sec:sanity}

Table~\ref{tab:sanity} reports Spearman correlations between original and perturbed attributions.
C2 vision crosses the historical decorrelation bar on all four tasks and language on three.
Repeated tokens are a possible explanation for StackCube's language result, but no control isolates this mechanism.
State is not shuffled, although its attribution gradients may change when another modality changes.
C1 correlations remain high under frozen-reference and backbone-cascade randomization.
These observations motivate input-only and gradient-only controls; they do not establish that model dependence is unnecessary or that input shuffling certifies grounding.

\begin{table}[tbp]
\centering
\footnotesize
\setlength{\tabcolsep}{4pt}
\begin{tabular}{@{}ll rrr@{}}
\toprule
\textbf{Task} & \textbf{Check} & $\boldsymbol{\rho}$\textbf{\,Vision} & $\boldsymbol{\rho}$\textbf{\,Lang.} & $\boldsymbol{\rho}$\textbf{\,State} \\
\midrule
PegInsertion & C1 & 0.953 & 0.950 & 0.786 \\
& C2 & \textbf{0.091} & \textbf{0.002} & 0.667 \\
\midrule
\rowcolor{gray!6}
PickCube & C1 & 0.937 & 0.903 & 0.750 \\
\rowcolor{gray!6}
& C2 & \textbf{0.160} & \textbf{0.087} & 0.810 \\
\midrule
PickYCB & C1 & 0.918 & 0.917 & 0.714 \\
& C2 & \textbf{0.166} & \textbf{0.141} & 0.762 \\
\midrule
\rowcolor{gray!6}
StackCube & C1 & 0.953 & 0.970 & 0.667 \\
\rowcolor{gray!6}
& C2 & \textbf{0.157} & 0.240 & 0.786 \\
\midrule
\midrule
\multicolumn{2}{@{}l}{C1 frozen (pooled)} & 0.929 & 0.862 & 0.643 \\
\multicolumn{2}{@{}l}{C1 cascade (pooled)} & 0.713 & 0.910 & 0.643 \\
\bottomrule
\end{tabular}
\caption{Sanity checks (Spearman $\rho$, where bold passes the $\leq$0.2 bar). The per-task rows are the standard C1 (final-layer re-init, moving target) and C2 (input shuffle) on action-norm-selected rows ($n$ = 67/62/85/75 per task). The bottom block holds the frozen-target and full-backbone cascade C1 variants, pooled over all four tasks and both seeds ($n{=}400$ each, no signal filter). Perturbed attributions are computed at $m{=}16$ (Section~\ref{sec:metrics}). The variant rows correlate against the regeneration run and carry the language-baseline and embedding-proxy deviations recorded in Section~\ref{sec:setup}. The released interval tables give a 95\% episode bootstrap interval of 0.906 to 0.940 for the pooled frozen vision median. The per-task C1 and C2 rows come from the pre-fix main pass and carry the same seed exposure disclosed in Section~\ref{sec:metrics}. The verification pass reproduces every standard-row verdict: C1 remains above the 0.2 bar at 0.916 to 0.965 on vision and language, C2 vision passes on all four tasks at 0.097 to 0.167, and C2 language passes on three of four, with StackCube borderline at 0.252.}
\label{tab:sanity}
\end{table}

\subsection{Qualitative Analysis}
\label{sec:qualitative}

Figure~\ref{fig:overlays} shows six representative vision overlays from a single checkpoint, the fine-tuned RDT-1B on PickCube at evaluation seed 242.
The top row is three successes and the bottom row is three failures, taken as the first three of each by episode index rather than hand-picked, with a representative mid-trajectory step shown per episode.
In the success panels, attribution concentrates on the target cube and the gripper's approach vector.
In the failure panels, the attribution pattern is mixed, with some panels scattering attribution across the scene and others concentrating on task objects despite the failed outcome.
Because both rows come from the same checkpoint and the same run, this is a descriptive within-checkpoint contrast rather than a cross-checkpoint one, and it shows that attribution placement alone does not separate success from failure here.

In inspected examples, the per-step language attribution bars peak on the content tokens that name the target object, such as \emph{red} and \emph{cube} on PickCube.
High correlation under backbone randomization cautions against treating token bars as independent evidence of learned grounding.
Spearman correlation does not estimate a fraction of model-dependent information.

A historical pipeline selected 16 rows from per-modality completeness residuals and named the groups vision misfocus, language misgrounding, state misreading, delayed credit, and success-well-grounded.
Those names are not validated behavioral diagnoses: an integration residual does not identify a perceptual, semantic, or temporal failure mechanism.
The groups should be interpreted only as numerical diagnostics.
The archived count of 3{,}881 generated images cannot be reconstructed from the released subset and is not used as scientific evidence.

\begin{figure}[tbp]
\centering
\includegraphics[width=0.92\textwidth,height=0.32\textheight,keepaspectratio]{figures/fig_b_overlays_grid.png}
\caption{Per-step vision IG overlays from a single checkpoint, the fine-tuned RDT-1B on PickCube at evaluation seed 242, where warm colors indicate high attribution. The top row is three successes and the bottom row is three failures, taken as the first three of each by episode index, with a representative mid-trajectory step per episode. Successful episodes often concentrate attribution on the target object and the gripper. The failure panels are mixed, some scattering attribution across the scene and some staying on task objects, so attribution placement alone does not separate success from failure in these examples. Because both rows come from the same checkpoint and run, this is a descriptive within-checkpoint contrast.}
\label{fig:overlays}
\end{figure}

\section{Response Geometry and Diagnostic Limits}
\label{sec:mechanisms}

The saved experiments support observations about attribution evaluation, but interpreting their causes requires separating score definitions, intervention effects and numerical validity.

\subsection{Action Displacement and the Quadratic Readout}
\label{sec:contraction}

Table~\ref{tab:faithfulness} reports smaller quadratic-score changes for the recorded 1B checkpoint than for the historical 170M cohort.
The checkpoints differ in size, training, action regime and trajectories, and only 170M uses action-norm filtering.
This does not identify the effect of competence or scale.

Direct displacement measurements show nonzero response to feature deletion on two tasks at both scales.
At $T=5$ and top-5\% vision deletion, historical median relative displacements range from 4.4 to 8.5\%.
Thus the near-zero score is not evidence that the action stays exactly fixed.
The identity $Q=-\lVert a-\mu\rVert^2/(2D\sigma^2)$ maps displacement to its numerical score.
Removing $1/D$ multiplies score changes by 512 without changing exact IG rankings, normalized AUC or relative completeness.
Operational interpretation requires stated action units, normalization and a justified measurement criterion.

Changing $T$ changes discretization over the complete denoising horizon; it does not add steps to one fixed trajectory.
Changing conditioning also differs from testing contraction under identical conditioning.
The endpoint medians from $T=2$ to $T=20$ do not show monotonic decay, but do not rule out denoiser contraction.
Episode-resampled 95\% intervals for changes of median relative displacement include appreciable decreases: PickCube-170M [$-23.82$, $31.04$]\%, PickCube-1B [$-10.11$, $37.88$]\%, StackCube-170M [$-30.10$, $73.97$]\%, and StackCube-1B [$-19.41$, $52.31$]\%.
The 1B curves pool two evaluation seeds, 170M uses one, and curves contain 8 to 40 episodes.
These are solver-resolution observations, not an identified comparison of denoiser mechanisms.

\begin{figure}[tbp]
\centering
\includegraphics[width=0.85\linewidth,height=0.38\textheight,keepaspectratio]{figures/fig_1b_vs_170m.png}
\caption{Historical quadratic-response AUCs on PickCube, with different populations as in Table~\ref{tab:faithfulness}. The 1B heights pool three evaluation seeds. Reference lines are historical vision bars without established chance calibration. The contrast changes checkpoint size, training, contexts and filtering and does not isolate competence.}
\label{fig:scale}
\end{figure}

\begin{figure}[tbp]
\centering
\includegraphics[width=0.85\linewidth,height=0.38\textheight,keepaspectratio]{figures/fig_displacement_vs_T.png}
\caption{Historical action displacement under top-5\% vision-token deletion across solver resolutions. Top: absolute displacement. Bottom: displacement divided by reference-action norm. Nonzero displacement shows changed actions, but varying resolution does not isolate contraction along a common trajectory. Bands are episode bootstrap intervals over 8 to 40 episodes per curve; 1B pools seeds 42 and 142.}
\label{fig:displacement}
\end{figure}

\subsection{Separating Attribution Target from Evaluation Response}
\label{sec:targets-results}

Table~\ref{tab:targets} retains historical native-score summaries.
Each target changes both the ranking and its evaluation response, so native AUC comparison does not isolate ranking quality.

\begin{table}[tbp]
\centering
\footnotesize
\setlength{\tabcolsep}{4pt}
\begin{tabular}{@{}l rrr rrrr@{}}
\toprule
 & \multicolumn{3}{c}{\textbf{Completeness}} & \multicolumn{4}{c}{\textbf{Faithfulness}} \\
\cmidrule(lr){2-4} \cmidrule(lr){5-8}
\textbf{Target} & med err & $\leq$3\% & $n$ & $\Delta F$@5 & \textbf{Ins} & \textbf{Del} & $n$ \\
\midrule
$Q$ & \textbf{1.86} & 76.8 & 699 & $-$0.0014 & \textbf{0.883} & 0.448 & 750 \\
\rowcolor{gray!6}
$\ell_2$ & 4.73 & 10.7 & 721 & $-$1.2299 & \textbf{0.755} & \textbf{0.284} & 774 \\
maxdev & 6.67 & 10.9 & 713 & $-$0.1562 & \textbf{0.723} & \textbf{0.273} & 765 \\
\bottomrule
\end{tabular}
\caption{Historical native-target summary on PickCube (170M, $m=64$, both seeds, vision). Magnitudes have target-specific units; bold denotes historical bars only. Completeness uses action-norm-selected rows; faithfulness includes all recorded rows. Alternative-target counts include resume duplicates, unlike the interval calculations. Reference-action norms differ on all 750 joined quadratic-versus-alternative calls, so exact matching is unestablished. These columns do not identify ranking advantage or an inherent accuracy-magnitude tradeoff.}
\label{tab:targets}
\end{table}

\paragraph{Fixed-ranking response transformation.}
Saved quadratic and $\ell_2$ perturbation curves can be rescored without recomputing IG or changing selected features.
For $D=512$, $\sigma^2=1$, and $\epsilon=10^{-12}$,
\begin{equation}
L=-\sqrt{-1024Q+\epsilon},\qquad Q=-\frac{L^2-\epsilon}{1024}.
\end{equation}
We normalize and integrate each transformed curve on its recorded grid, retaining the last record per seed, episode and policy-call key.
Each arm contains 750 distinct calls.
For quadratic-target rankings, median deletion AUC changes from 0.447903 under $Q$ to 0.290759 under $L$; insertion changes from 0.883442 to 0.755679.
For $\ell_2$ rankings, deletion changes from 0.283007 under $L$ to 0.440376 under $Q$; insertion changes from 0.756490 to 0.884070.
The historical deletion verdict reverses in both directions solely through response transformation.
This requires no cross-arm context match and preserves the recorded curves' provenance limitations.
It does not establish ranking identity, equivalence or efficacy.

\paragraph{Numerical and scientific interpretation.}
The historical alternative-target cohorts have larger relative completeness residuals at $m=64$.
Stabilized $\ell_2$ is smooth, so an actual square-root kink does not explain its error.
High curvature, quadrature, accumulation and model-forward precision, and cohort differences remain unseparated.
Scaling a score multiplies exact IG and score changes while preserving ranking and relative completeness; numerical magnitude alone is not recovered information.
A valid comparison must evaluate quadratic-IG and $\ell_2$-IG on authenticated paired contexts and noise under a common response, alongside random and simple rankings.
Positive and negative outcomes are both informative about attribution beyond response geometry.

\subsection{What Randomization Correlations Establish}
\label{sec:anchoring}

Saved frozen-reference and backbone-cascade variants retain high correlations.
Freezing the reference addresses one moving-target explanation but does not isolate every source of correlation.
Encoders and adaptors remain learned, integration budgets differ, and overwritten sidecars limit context identity.
The values in Table~\ref{tab:sanity} are descriptive within these limits.

\begin{figure}[tbp]
\centering
\includegraphics[width=0.85\linewidth,height=0.38\textheight,keepaspectratio]{figures/fig_c1_variants.png}
\caption{Historical model-randomization correlations. Frozen-reference and cascade values use saved records; standard C1 and C2 preserve unavailable original summaries. Variants change randomized parameters and may differ in source contexts or budgets. Persistent correlation does not uniquely identify input anchoring; low C2 correlation does not establish learned-model dependence.}
\label{fig:c1}
\end{figure}

The factor $(x_i-x_i')$ in Eq.~\ref{eq:ig} permits shared structure but does not force a particular ranking.
Integrated gradients can dominate rankings, and the input features depend on retained learned encoders and adaptors.
No direct input-only or gradient-only comparison establishes dominance of the multiplier.
Nor is Spearman correlation a fraction of model-independent information.

Input shuffling asks whether a map changes when input arrangement changes.
A map defined only from input magnitudes can pass while ignoring learned policy parameters.
C2 therefore cannot replace model-dependence checks or certify grounding.

Table~\ref{tab:baseline} shows positive but imperfect black-gray association and greater differences involving blur.
With only two episode clusters per task, it does not establish equivalence.
Baseline changes affect the interpolation path and gradients as well as the input difference, so the mechanism is unisolated.

\begin{table}[tbp]
\centering
\small
\begin{tabular}{@{}l rrr@{}}
\toprule
\textbf{Task} & \textbf{black vs gray} & \textbf{black vs blur} & \textbf{gray vs blur} \\
\midrule
PickCube & 0.825 & 0.539 & 0.455 \\
\rowcolor{gray!6}
StackCube & 0.814 & 0.546 & 0.496 \\
\bottomrule
\end{tabular}
\caption{Historical vision baseline sensitivity: median patch-ranking Spearman correlation over the first 30 recorded calls per task at $m=64$. Pre-fix sidecars leave source seed identity uncertain. Within a row, all baselines use one loaded context and noise seed. Black-gray agreement is imperfect, and two episode clusters per task do not establish equivalence or unchanged top-$k$ selections.}
\label{tab:baseline}
\end{table}

\subsection{Limits of the Historical Integration-Budget Comparison}
\label{sec:budget}

Table~\ref{tab:m128} preserves historical completeness pass-rate summaries.
The rows previously called matched truncate $m=128$ to four calls but do not pair identical observations, episodes, language baselines or runtimes.
The $m=64$ comparator is the larger main pass; early $m=128$ uses eight episodes per seed and its original PickCube/StackCube records are missing.
Differences are compatible with a budget effect but cannot identify it.
Full-view rates also change the trajectory-depth distribution.
Later saved $m=128$ cohorts cannot recover the absent early records.
\begin{table}[tbp]
\centering
\footnotesize
\setlength{\tabcolsep}{4.5pt}
\begin{tabular}{@{}ll rrr@{}}
\toprule
\textbf{Task} & \textbf{View} & $\boldsymbol{m{=}64}$ & $\boldsymbol{m{=}128}$ & $\boldsymbol{n_{128}}$ \\
\midrule
PickCube & matched & 68.6 & \textbf{80.5} & 41 \\
StackCube & matched & 59.9 & \textbf{89.8} & 49 \\
\midrule
PickCube & full & 68.6 & 72.9 & 377 \\
StackCube & full & 59.9 & 62.0 & 379 \\
PegInsertion & full & 75.8 & 64.0\,/\,77.0 & 189\,/\,178 \\
PickYCB & full & 69.8 & 73.7\,/\,71.8 & 194\,/\,195 \\
\bottomrule
\end{tabular}
\caption{Vision completeness pass rate (\% of action-norm-selected rows at or under 3\% error) at $m{=}64$ and $m{=}128$. The ``matched'' view truncates $m{=}128$ episodes to the $m{=}64$ episode shape (first four policy calls), which aligns maximum call depth but does not pair contexts, and bold marks its two cells. The PickCube and StackCube cells in both the matched and full views, together with their $m{=}128$ faithfulness AUCs quoted in the text, come from the decommissioned main-pass $m{=}128$ subset and are not regenerable from the released records. The released records cover the PegInsertion and PickYCB $m{=}128$ runs only, as the reproducibility statement records. The ``full'' view keeps every call of each $m{=}128$ episode, a row mix containing deep-trajectory steps absent from the $m{=}64$ pass. PegInsertion and PickYCB are reported per seed (42\,/\,142). The PickYCB $m{=}128$ run reused the PickCube instruction embedding (Section~\ref{sec:setup}), and no language metrics are tabulated from it. The PegInsertion counts retain the resume-replay duplicate rows documented in the released data README, and removing them shifts the seed 142 PegInsertion cell from 77.0 on $n{=}178$ to 76.2 on $n{=}160$ with no verdict change.}
\label{tab:m128}
\end{table}

These are descriptive cohort comparisons, not evidence of smooth path integrands or a budget-only cause.
A convergence study requires identical inputs, baselines, checkpoint and stored noise, with separate control of forward precision, accumulation and quadrature.
Coordinate, ranking and response stability checks must supplement scalar completeness.

%% ================================================================
\FloatBarrier
\section{Discussion}

Evaluation should name the attribution target, ranking, intervention, response, normalization and population separately.
Exact rescoring shows that normalized AUC depends on response geometry as well as ranking.
Our counterexample shows why fixed bars cannot substitute for matched controls.
Randomization results also need input-only and gradient-only comparisons before persistent structure is assigned a dominant cause.
These requirements support informative comparisons of attribution methods; historical threshold crossings alone do not settle them.

\paragraph{Limitations.}
The cross-scale displacement comparison covers two tasks at two seeds at 1B scale, and the target ablation covers one task.
Both invite replication at scale.
The 1B evaluation seeds vary the rollouts and the sampling on one fixed checkpoint, so they bound rollout and sampling variance and not generalization across training initializations.
The tables report medians or pooled pass rates, and rows within an episode are correlated, so row-level $n$ overstates the effective sample size.
The 1{,}199 pooled 1B rows, for example, come from 60 episodes.
We mitigate this by reporting per-seed values where populations differ and by scoping claims to directions and orders of magnitude.
The tables also make many pass or fail comparisons against pre-specified bars without a multiplicity correction, and we read them as descriptive summaries rather than inferential tests.
Episode-level bootstrap intervals accompany Figure~\ref{fig:displacement} and ship as interval tables for every released run that backs a published table, while the decommissioned runs remain without intervals.
The distribution figures (Figures~\ref{fig:completeness} and~\ref{fig:auc}) and the in-text $m{=}128$ AUC values are not accompanied by interval tables.
The historical pretrained-only 170M has very low recorded success, limiting behavioral interpretation.
State completeness misses its median bar and has no perturbation AUC evaluation.
Eight active state dimensions permit a nine-point deletion grid, so omission is a scope choice, not an impossibility.
C2 changes vision and language only, with language shuffled after encoding; modality differences remain unexplained.
Baseline sensitivity is vision-only with few independent episodes.
Follow-up language results span the baseline and proxy-embedding deviations in Section~\ref{sec:setup}.
Numerical diagnostic groups are not validated behavioral failure labels.
Historical maps lack same-context high-precision checks, and missing sidecars, checkpoints and environment identities prevent exact reproduction of several campaigns.
All results concern one model family in ManiSkill3 simulation; 1B faithfulness covers one task.
\paragraph{Engineering notes.}
The historical execution made the following engineering choices.
The language-baseline preference arose from preliminary tests whose records are unavailable; its efficacy is not independently verified.
Fixed noise defines the sample-conditioned estimand; attribution of a noise expectation is a different possible construction.
The loop attributes each policy call and executes at most 16 subsampled steps before replanning or termination.
Historical timings need a current benchmark before selecting new numerical settings or sample sizes.

\paragraph{Future directions.}
A paired study of quadratic-IG, $\ell_2$-IG, random order, input-only magnitude and path-gradient-only rankings (omitting the IG input multiplier) would separate ranking information from response geometry and anchoring.
Sampling should target independent episodes, a meaningful effect and uncertainty bound, with explicit failure accounting.
Smoother model constructions~\citep{jha2021smoother,jha2022shaping} and adaptive integration~\citep{walker2024idg} suggest numerical studies, but their benefits for this sampler and score require measurement.
Broader tasks, checkpoints and real-robot evaluation would test transfer.

%% ================================================================
\section{Conclusion}
We apply per-policy-call, per-modality IG to a fixed-noise diffusion sampler and examine what evaluation measures.
Changing only the response of recorded perturbation curves reverses historical normalized-AUC verdicts.
An analytic example shows that random ranking need not yield AUC 0.5.
Displacement and randomization observations motivate controls but do not establish exclusive denoiser or anchoring mechanisms.
Efficacy requires matched rankings and controls under a common response, numerical validation and authenticated contexts.
Saved-record analyses preserve useful evidence while making these dependencies explicit.

\section*{Ethical Statement}

This work analyzes pretrained robotic manipulation policies in simulation.
No human subjects were involved.
The interpretability methods presented here are intended to improve the transparency and safety of robotic systems.

\section*{Reproducibility Statement}

Per-step completeness metrics are emitted as JSONL records with one row per policy call, carrying the per-modality errors, IG sums, action norm, and wall time.
The faithfulness and sanity post-processing stages write companion JSONL files keyed to the same steps.
The original protocol asked for one consolidated metrics file, and this multi-file layout is a deviation from that requirement.
The published table values apply the action-norm-selected filter of Section~\ref{sec:setup}, and the matched-view and ablation columns derive from their records under the filters stated in their captions.
Episode seeds are fixed at 42 and 142, with 242 added for the third RDT-1B evaluation seed, where episode $e$ uses seed $42{+}e$, $142{+}e$, or $242{+}e$, and the diffusion noise is seeded per forward pass.
The repository preserves scripts for the historical full pass, faithfulness, sanity, target-ablation and displacement stages. A working command and fixed seed do not guarantee reproduction without the corresponding checkpoint, conditioning artifacts and runtime identity.
The seed 142 and 242 RDT-1B faithfulness runs, the nominally matched $Q$ ablation runs, and the seed 142 displacement runs used the same entry points and flags at the seeds and scopes recorded in the released data README.
Historical 1B execution required the authors' fine-tuned ManiSkill checkpoint; base-weight or local LoRA fallbacks do not reproduce it.
Current support requires strict checkpoint identity and verified loading, gradient and rendering behavior; fixed seeds cannot recover missing weights or patches.
The committed records are not a byte-level reproduction of every table.
Six artifact families were produced on cloud GPUs that have since been decommissioned.
They are the main evaluation pass and its faithfulness records, the original $m{=}128$ PickCube and StackCube subset, the standard sanity runs, the initial RDT-1B scale-check run, the per-step attribution sidecars, and the overlay images.
None of them are part of the release.
The records we release are the regeneration run of Section~\ref{sec:setup}, which reproduces the target ablation and corroborates the main-pass protocol at a smaller episode budget.
The release also holds the RDT-1B faithfulness records at all three evaluation seeds, the $m{=}128$ PegInsertionSide and PickSingleYCB runs, the frozen-target and cascade sanity runs, and the baseline-sensitivity and target-ablation runs including the nominally matched $Q$ pair.
It further holds the six displacement records behind Figure~\ref{fig:displacement}, the RDT-1B base attribution runs behind the second displacement seed, the Month 5 1B reproduction records, and a validation record for the displacement identity.
The RDT-170M main-pass values were computed from the now-decommissioned main-pass records and are not regenerable from the committed data.
These are the completeness values in Table~\ref{tab:correctness}, the RDT-170M block of Table~\ref{tab:faithfulness}, and the standard C1 and C2 rows of Table~\ref{tab:sanity}.
The three-seed RDT-1B rows of Table~\ref{tab:faithfulness} and the frozen-target and cascade rows of Table~\ref{tab:sanity} are computed from the committed records and regenerate exactly from them, with bootstrap intervals in the released interval tables.
The released analysis notebooks preserve the Table~\ref{tab:correctness} values exactly and retain the other analyses over broader unfiltered row populations.
A README maps each committed file to the run it represents.
Two committed records are retained with a known defect rather than silently replaced: the seed-42 $m{=}128$ faithfulness files consumed seed-142 sidecars through the pre-fix shared filenames, as indicated by the released audit script's incompatible gap fingerprint of Section~\ref{sec:metrics}, and the README marks them accordingly; no published value derives from them.
The separate Month 7 campaign used seed-tagged sidecars and agrees with many historical descriptions, but changes PickCube's borderline language-deletion verdict.
Missing sidecars are not authenticated or recovered by fixed seeds alone.
Arithmetic reproduction from records is distinct from model reproduction under pinned code, checkpoints and exact inputs.
\ifanonymous
The identified public repository link is omitted from this anonymous draft. The review artifact must be assembled and audited before submission.
\else
Code, records, interval tables, and the analysis notebooks are available at \url{https://github.com/itscool2b/the-readout-not-the-denoiser-repo}.
\fi

\ifanonymous\else
\section*{Acknowledgments}

I am sincerely grateful to Professor Sumit K.\ Jha, Department of Computer and Information Science and Engineering, University of Florida, and to Brian Rodriguez, University of Florida, for mentoring and guiding me throughout this seven-month project. I thank them for their time and technical advice.
\fi

%% ================================================================
\bibliographystyle{tmlr}
\bibliography{references}

\end{document}
